\documentclass[11pt,a4paper]{article}

% ── Packages ──────────────────────────────────────────────────
\usepackage[utf8]{inputenc}
\usepackage[T1]{fontenc}
\usepackage[margin=2.5cm]{geometry}
\usepackage{amsmath, amssymb, amsfonts}
\usepackage{booktabs}         % nicer tables
\usepackage{enumitem}         % list customisation
\usepackage{hyperref}
\usepackage{xcolor}
\usepackage{titlesec}
\usepackage{parskip}          % paragraph spacing instead of indent
\usepackage{float}            % precise float placement
% microtypographic refinements skipped (font compatibility)

% ── Colours & Hyperlinks ─────────────────────────────────────
\definecolor{secblue}{HTML}{1a3c6e}
\hypersetup{
    colorlinks=true,
    linkcolor=secblue,
    urlcolor=secblue,
    citecolor=secblue,
}

% ── Section formatting ───────────────────────────────────────
\titleformat{\section}
  {\Large\bfseries\color{secblue}}{\thesection}{1em}{}
\titleformat{\subsection}
  {\large\bfseries\color{secblue!80!black}}{\thesubsection}{1em}{}
\titleformat{\subsubsection}
  {\normalsize\bfseries\color{secblue!60!black}}{\thesubsubsection}{1em}{}

% ── Custom commands ──────────────────────────────────────────
\newcommand{\cm}{\mathrm{cm}}
\newcommand{\contact}{\mathrm{contact}}
\newcommand{\grav}{\mathrm{gravity}}
\newcommand{\diss}{\mathrm{dissipation}}
\newcommand{\leading}{\mathrm{leading}}
\newcommand{\trailing}{\mathrm{trailing}}

% ── Title ────────────────────────────────────────────────────
\title{%
  \vspace{-1cm}
  {\LARGE\bfseries IYPT 2026 --- Problem \#6: Flipo Flip}\\[6pt]
  {\large Theoretical Framework \& Simulation Model}
}
\author{}
\date{}

\begin{document}
\maketitle
\thispagestyle{empty}

% ═════════════════════════════════════════════════════════════
\section{Problem Statement}
% ═════════════════════════════════════════════════════════════

A Flipo Flip toy can roll for multiple turns even though its shape is not circular.
Investigate how its motion depends on parameters such as geometry and the initial
release conditions.

% ═════════════════════════════════════════════════════════════
\section{Geometric Model}
% ═════════════════════════════════════════════════════════════

We model the Flipo Flip as a square prism (rectangular cross-section) with cylindrical
caps.  Specifically, the toy's curved surfaces are the two ``outer shells'' of a
Steinmetz solid---the intersection region of two orthogonal cylinders of radius~$R$.
These curved caps are attached to opposite faces of the prism.

\subsection{2D Cross-Section}

We work entirely in 2D\@.  In the plane perpendicular to the rolling axis, the toy
appears as:
\begin{itemize}[nosep]
  \item A rectangle of width~$w$ (the prism's face along the rolling direction) and
        height~$h$ (the prism's face perpendicular to it), where $w \neq h$ in general.
  \item Two circular arcs of radius~$R$, each capping one of the short ($w$-wide) edges
        of the rectangle.
\end{itemize}

\textbf{Note:} The prism is \emph{not} generally inscribed in the cylinder---that would
be a special case.  In general, two edges of the rectangle lie on the circle of
radius~$R$, while the other two are interior to it.  The distance from the centre of the
prism to the cylinder's axis (the ``apothem'') is determined by the prism
side length and~$R$.

\subsection{Moment of Inertia (Simplifying Approximation)}

For simplicity, we approximate the moment of inertia as that of the rectangular prism
alone, ignoring the mass of the thin curved caps.  This gives us clean, closed-form
expressions for $I_{\cm}$ of a uniform rectangular solid, which we then shift to
whichever contact point or axis is relevant using the parallel axis theorem.  The curved
caps still define the rolling geometry (radius~$R$), but their contribution to inertia is
neglected.

% ═════════════════════════════════════════════════════════════
\section{Phases of Motion}
% ═════════════════════════════════════════════════════════════

The rolling motion naturally decomposes into two alternating phases.

% ─────────────────────────────────────────────────────────────
\subsection{Phase 1: Pure Patch Rolling}
% ─────────────────────────────────────────────────────────────

During this phase, one of the two curved (cylindrical) surfaces is in contact with the
ground and rolls without slipping.  The other curved surface is in the air and contributes
only through its inertia.

\begin{description}[style=nextline, leftmargin=1.5cm]
  \item[Contact geometry]
    The contact point traces along the circular arc of radius~$R$.
  \item[Instantaneous axis of rotation]
    A line parallel to the contact line, passing through the centre of the parent cylinder
    (i.e.\ the centre of the circle of radius~$R$ from which this curved surface
    originates).
  \item[Dynamics]
    Standard rolling-without-slipping of a circular profile, with the complication that the
    centre of mass does not coincide with the geometric centre of the rolling circle.
\end{description}

This phase continues until the flat bottom face of the prism reaches the ground---i.e.\
until the angle between the prism face and the floor reaches zero.

% ─────────────────────────────────────────────────────────────
\subsection{Phase 2: Patch Switch}
% ─────────────────────────────────────────────────────────────

When the curved surface has rolled far enough that the flat prism face becomes tangent
to the ground, the motion transitions to a rotation about the edges of the rectangle:

\begin{enumerate}[nosep]
  \item \textbf{Trailing-edge pivot:} The toy pivots about the \emph{trailing edge}---the edge at
        the rear, where the just-finished arc meets the prism face.  The opposite curved
        surface swings forward and down.
  \item \textbf{Leading-edge pivot:} When the full width~$w$ of the prism face contacts the
        floor, the pivot instantaneously transfers to the \emph{leading edge}---the edge
        on the far (forward) side of the flat face, where the next arc begins.  This new
        pivot point is in front of the previous one in the direction of travel.
\end{enumerate}

These two sub-phases are not equal in duration; their relative timing depends on the
geometry and the angular velocity at entry.

The leading-edge rotation continues until the new curved surface engages with the ground
and Pure Patch Rolling resumes.

\begin{description}[style=nextline, leftmargin=1.5cm]
  \item[Instantaneous axis of rotation]
    First at the trailing edge, then switches to the leading edge (forward) at the moment
    the full flat face is flush with the floor.
  \item[No-slip condition]
    Must hold at the pivot edge throughout.
  \item[Critical angle]
    The transition angle---from Patch Switch back to Pure Patch Rolling---is set by the
    geometry.  It is the angle at which the next circular arc first becomes tangent to the
    floor, which depends on~$R$, $w$, and~$h$.
\end{description}

% ═════════════════════════════════════════════════════════════
\section{Theoretical Framework}
% ═════════════════════════════════════════════════════════════

% ─────────────────────────────────────────────────────────────
\subsection{Equations of Motion}
% ─────────────────────────────────────────────────────────────

For each phase, we write the rotational equation of motion about the instantaneous
contact point/axis:
\begin{equation}
  I_{\contact}\,\dot{\omega}
    = \tau_{\grav} - \tau_{\diss}
\end{equation}
where $I_{\contact}$ is the moment of inertia about the current contact point (computed
via the parallel axis theorem from the centre-of-mass value), and $\tau_{\grav}$ is the
gravitational torque about that point.

% ─────────────────────────────────────────────────────────────
\subsection{Energy Dissipation}
% ─────────────────────────────────────────────────────────────

Energy is lost through four mechanisms:

\begin{enumerate}
  \item \textbf{Rolling friction} during Pure Patch Rolling, modelled as a rolling
        resistance coefficient~$\mu_r$.

  \item \textbf{Pivot friction} during Patch Switch.  Although the contact point is fixed
        (an edge, not a rolling arc), the edge is not infinitely sharp: it has an effective
        radius of curvature~$r_{\mathrm{edge}}$, and presses into the surface under the
        normal force~$N$.  The resulting asymmetric deformation at the contact produces a
        resistive torque
        \begin{equation}
          \tau_{\mathrm{pivot}} = \mu_p \, N \, r_{\mathrm{edge}}
        \end{equation}
        where $\mu_p$ is a pivot friction coefficient.  The physical mechanism is the same
        as for rolling friction (asymmetric contact deformation) but at the smaller length
        scale of the edge.  In practice, Flipo Flip toys have visibly rounded edges, so
        $r_{\mathrm{edge}}$ is a real, measurable geometric parameter---not an abstraction.
        Combined with $\mu_p$, this makes pivot friction a physically grounded and
        potentially significant contribution, especially on softer surfaces.

        As with the Pure Patch Rolling phase, the no-slip condition must also be checked at
        the pivot edge: the simulation computes the tangential force required to prevent
        slipping and compares it to~$\mu_s N$.  If violated, slip is detected and the
        physical simulation halts.

  \item \textbf{Impulsive losses at the axis switch}---when the pivot transfers from the
        trailing edge to the leading edge during Patch Switch, we conserve angular momentum
        about the new contact point.  That is, $\omega_{\mathrm{after}}$ is determined by:
        \begin{equation}
          I_{\leading}\cdot\omega_{\mathrm{after}}
            = L_{\leading}^{\mathrm{before}}
        \end{equation}
        where $L_{\leading}^{\mathrm{before}}$ is the angular momentum about the leading
        edge computed from the pre-impact state (rotating about the trailing edge).
        Because $I_{\leading}\neq I_{\trailing}$ and the velocity field changes, this
        generically implies a kinetic energy loss that is \emph{entirely predicted by
        geometry}---not a free parameter.

  \item \textbf{Additional impact dissipation} $e\in[0,1]$---a small correction factor
        applied on top of the angular momentum conservation result, accounting for energy
        lost to vibrations and sound during the impact.  (Material deformation is
        negligible for the hard plastic of typical Flipo Flip toys.)  With
        $e=1$, no extra loss is applied beyond what angular momentum conservation already
        implies.
\end{enumerate}

% ─────────────────────────────────────────────────────────────
\subsection{No-Slip Condition and Slip Detection}
% ─────────────────────────────────────────────────────────────

Throughout all phases, we assume no-slip at the contact point.  At each time step, the
simulation computes the required friction force~$f$ to maintain this constraint, and
compares it to the maximum available static friction~$\mu_s N$ (where $N$ is the normal
force at the contact).

\begin{itemize}[nosep]
  \item If $f \leq \mu_s N$: the no-slip model holds and integration continues.
  \item If $f > \mu_s N$: slip is detected.  The physical simulation halts, because our
        rolling model cannot describe the post-slip dynamics accurately.  For
        visualisation purposes only, a simplified ``fall'' animation plays out (e.g.\ the
        toy sliding and toppling under gravity with some damping), so the user can see the
        toy come to rest rather than freezing mid-frame.  This visual-only phase makes no
        physical predictions.
\end{itemize}

% ─────────────────────────────────────────────────────────────
\subsection{Matching Conditions Between Phases}
% ─────────────────────────────────────────────────────────────

At each phase transition, we must match:
\begin{itemize}[nosep]
  \item The angular velocity (converting between the two different rotation axes using
        conservation of angular momentum, with possible impulsive losses).
  \item The angular position (the critical angle at which the transition occurs, determined
        purely by geometry).
\end{itemize}

% ═════════════════════════════════════════════════════════════
\section{Parameters}
% ═════════════════════════════════════════════════════════════

\begin{table}[H]
\centering
\renewcommand{\arraystretch}{1.3}
\begin{tabular}{@{} c l @{}}
\toprule
\textbf{Symbol} & \textbf{Meaning} \\
\midrule
$R$        & Radius of the parent cylinder (determines curvature of the caps) \\
$w$        & Width of the prism face along the rolling direction \\
$h$        & Height of the prism face perpendicular to rolling \\
$\rho$     & Material density (or equivalently, total mass~$m$) \\
$\omega_0$ & Initial angular velocity at release \\
$\mu_r$    & Rolling friction coefficient \\
$\mu_p$    & Pivot friction coefficient (at edges during Patch Switch) \\
$r_{\mathrm{edge}}$ & Effective radius of curvature of the prism edges \\
$\mu_s$    & Static friction coefficient (used for slip detection) \\
$e$        & Extra energy retention factor at each Patch Switch impact \\
           & \quad ($e=1$: no extra loss; $e<1$: additional loss to vibrations, etc.) \\
\bottomrule
\end{tabular}
\caption{Summary of model parameters.}
\label{tab:parameters}
\end{table}

% ═════════════════════════════════════════════════════════════
\section{Initial Condition}
% ═════════════════════════════════════════════════════════════

The simulation begins with the toy standing vertically (balanced on one curved cap),
given an initial angular velocity~$\omega_0$ that tips it into its rolling motion.
Experimentally, this velocity is imparted by accelerating the toy as it is brought to
horizontal, but for the simulation we simply prescribe~$\omega_0$ at~$t=0$.

% ═════════════════════════════════════════════════════════════
\section{Goal: Interactive Simulation}
% ═════════════════════════════════════════════════════════════

This model is intended to be complete enough to drive a real-time, interactive simulation
of the Flipo Flip's motion.  The simulation should:

\begin{enumerate}
  \item Numerically integrate the differential equations of motion through alternating
        Pure Patch Rolling and Patch Switch phases, detecting phase transitions by
        geometric criteria.
  \item Visualise the 2D cross-section rolling across a flat surface, showing the rectangle
        and arcs in motion.
  \item Let the user adjust parameters ($R$, $w$, $h$, $\omega_0$, $\mu_r$, $\mu_p$,
        $\mu_s$, $e$) with sliders and immediately see the effect on the motion.
  \item Display diagnostic outputs: number of complete rolls, angular velocity over time,
        energy over time (showing the stepwise losses at each Patch Switch).
  \item Detect and report slip: when the required friction exceeds~$\mu_s N$, halt the
        physical simulation, display a warning, and play a visual-only fall/slide animation
        so the toy doesn't freeze in mid-air.
\end{enumerate}

% ═════════════════════════════════════════════════════════════
\section{Key Predictions the Model Should Capture}
% ═════════════════════════════════════════════════════════════

\begin{enumerate}
  \item Number of complete rolls as a function of~$\omega_0$ and geometry.
  \item Time evolution of angular velocity, showing the stepwise energy loss at each
        Patch Switch.
  \item Sensitivity to geometry---how the ratio~$w/R$ and~$h/R$ affect rolling efficiency.
  \item Dominant dissipation mechanism---whether rolling friction, pivot friction, or
        impact losses matter more.
\end{enumerate}

\end{document}
